\section{GAS 调试面板代码学习路径}

\subsection{推荐阅读顺序}
\begin{itemize}
	\item 先看入口：是谁创建面板、谁控制显示
	\item 再看核心：面板怎么拿 \texttt{ASC} 和 \texttt{AttributeSet}
	\item 再看刷新：数据怎么更新到 UI
	\item 最后看扩展点：以后自己往里加什么最方便
\end{itemize}

按这个顺序阅读，会比较容易建立整体结构，而不是一上来陷进某个函数细节里。

\subsection{第一步：先看是谁打开这个面板的}
先看文件：

\begin{itemize}
	\item \texttt{Source/Aura/Private/Player/AuraPlayerController.cpp}
\end{itemize}

重点看：
\begin{itemize}
	\item \texttt{AAuraPlayerController::ToggleAttributeDebugPanel()}
	\item \texttt{SetupInputComponent()} 里绑定按键的地方
\end{itemize}

这一层主要回答两个问题：
\begin{enumerate}
	\item 为什么按下 \texttt{F12} 会生效？
	\item 按键事件最后是如何转发给 \texttt{HUD} 的？
\end{enumerate}

可以把这一层理解成：

\begin{itemize}
	\item \texttt{PlayerController} 负责输入
	\item 输入触发后，调用 \texttt{HUD} 去显示或隐藏调试面板
\end{itemize}

\subsection{第二步：再看 HUD 是怎么创建这个面板的}
然后看：

\begin{itemize}
	\item \texttt{Source/Aura/Public/UI/HUD/AuraHUD.h}
	\item \texttt{Source/Aura/Private/UI/HUD/AuraHUD.cpp}
\end{itemize}

重点看两个函数：

\begin{itemize}
	\item \texttt{AAuraHUD::InitOverlay(...)}
	\item \texttt{AAuraHUD::InitAttributeDebugPanel(...)}
\end{itemize}

这里是整个调用链里最关键的一步。你会看到：

\begin{enumerate}
	\item HUD 先初始化正常的 Overlay
	\item 然后顺手创建 \texttt{UAuraAttributeDebugWidget}
	\item 再把 \texttt{ASC} 和 \texttt{AttributeSet} 传进去
\end{enumerate}

这一层可以概括为：

\begin{itemize}
	\item \texttt{PlayerController} 负责“开关”
	\item \texttt{HUD} 负责“创建和持有面板”
\end{itemize}

\subsection{第三步：看 Widget 本体}
然后重点看：
\begin{marker}
	\begin{itemize}
		\item \texttt{Source/Aura/Public/UI/Widget/AuraAttributeDebugWidget.h}
		\item \texttt{Source/Aura/Private/UI/Widget/AuraAttributeDebugWidget.cpp}
	\end{itemize}
\end{marker}
这是整个功能的核心部分。

建议按下面顺序读函数。
\begin{enumerate}
	
	\item \textbf{生命周期函数}\\
	先看：
	
	\begin{itemize}
		\item \texttt{InitializeWithAbilitySystem(...)}
		\item \texttt{NativeOnInitialized()}
		\item \texttt{NativeTick(...)}
		\item \texttt{NativeDestruct()}
	\end{itemize}
	
	你要先搞懂它的“活法”：
	
	\begin{itemize}
		\item 什么时候创建
		\item 什么时候绑定 \texttt{ASC}
		\item 什么时候定时刷新
		\item 什么时候解绑委托
	\end{itemize}
	
	如果这组函数读明白了，后面很多刷新逻辑都会更容易理解。
	
	\item \textbf{UI 是怎么纯 C++ 搭出来的}
	重点看：
	
	\begin{itemize}
		\item \texttt{BuildWidgetTree()}
	\end{itemize}
	
	这个函数是在不依赖蓝图控件树的情况下，直接在 C++ 里创建 UI 结构。这里会创建：
	
	\begin{itemize}
		\item \texttt{CanvasPanel}
		\item \texttt{Border}
		\item \texttt{ScrollBox}
		\item \texttt{VerticalBox}
		\item 多个 \texttt{TextBlock}
	\end{itemize}
	
	如果把这个函数看懂，你对 UE 的 UMG 底层组织方式会清晰很多。
	
	建议重点关注这些对象之间的关系：
	
	\begin{itemize}
		\item \texttt{RootWidget}
		\item \texttt{PanelBorder}
		\item \texttt{DebugScrollBox}
		\item \texttt{DebugContentBox}
		\item 各种 \texttt{TextBlock}
	\end{itemize}
	
	本质上，这里就是在手动搭一棵 Widget Tree。
	
	\item \textbf{总刷新入口}
	重点看：
	
	\begin{itemize}
		\item \texttt{RefreshValues()}
	\end{itemize}
	
	这个函数可以理解为总控刷新器。它会统一去做：
	
	\begin{itemize}
		\item 刷新属性
		\item 刷新 Tag
		\item 刷新 Gameplay Effect
		\item 刷新 Ability
	\end{itemize}
	
	读到这里时，建议在脑子里建立这样一条总链路：
	
	\begin{center}
		\texttt{ASC / AttributeSet} $\rightarrow$ 读取数据 $\rightarrow$ 写入缓存 $\rightarrow$ \texttt{UpdateDisplayedValues()} $\rightarrow$ 更新到 UI
	\end{center}
	
	\item \textbf{每个分区的数据怎么来的}
	依次看：
	
	\begin{itemize}
		\item \texttt{RefreshOwnedTags()}
		\item \texttt{RefreshActiveGameplayEffects()}
		\item \texttt{RefreshActivatableAbilities()}
	\end{itemize}
	
	这三个函数分别对应三个调试区块。
	
	\paragraph{1. \texttt{RefreshOwnedTags()}}
	这个函数主要看：
	
	\begin{itemize}
		\item 如何从 \texttt{AbilitySystemComponent} 里读取 \texttt{OwnedGameplayTags}
	\end{itemize}
	
	\paragraph{2. \texttt{RefreshActiveGameplayEffects()}}
	这个函数值得重点学。你会看到如何读取当前所有 Active Gameplay Effects，以及如何提取：
	
	\begin{itemize}
		\item \texttt{Stack}
		\item \texttt{Duration}
		\item \texttt{RemainingTime}
		\item \texttt{Period}
		\item \texttt{GrantedTags}
	\end{itemize}
	
	这部分对理解 GAS 的运行态信息很有帮助。
	
	\paragraph{3. \texttt{RefreshActivatableAbilities()}}
	这个函数主要学习：
	
	\begin{itemize}
		\item 如何枚举 \texttt{GetActivatableAbilities()}
		\item 如何读取 Ability 的名称、等级、输入 ID、激活状态
	\end{itemize}
	
	\item \textbf{最终怎么显示到屏幕}
	最后看：
	
	\begin{itemize}
		\item \texttt{UpdateDisplayedValues()}
	\end{itemize}
	
	这个函数负责把前面缓存好的内容真正塞到 UI 控件里，例如：
	
	\begin{itemize}
		\item \texttt{HealthText}
		\item \texttt{ManaText}
		\item \texttt{TagValuesText}
		\item \texttt{EffectValuesText}
		\item \texttt{AbilityValuesText}
	\end{itemize}
	
	如果你以后要扩展新字段，一般就是按这个套路：
	
	\begin{enumerate}
		\item 在头文件里新增缓存变量或控件指针
		\item 在某个 \texttt{RefreshXXX()} 函数里填充数据
		\item 在 \texttt{UpdateDisplayedValues()} 里把它显示出来
	\end{enumerate}
	
\end{enumerate}

\subsection{第四步：最后看属性变化的实时绑定}
看下面这些函数：

\begin{itemize}
	\item \texttt{BindAttributeDelegates()}
	\item \texttt{UnbindAttributeDelegates()}
	\item \texttt{HandleHealthChanged(...)}
	\item \texttt{HandleMaxHealthChanged(...)}
	\item \texttt{HandleManaChanged(...)}
	\item \texttt{HandleMaxManaChanged(...)}
\end{itemize}

这里的重点是理解：

\subsubsection{为什么 Attribute 用 Delegate}
因为属性变化频繁，而且 GAS 已经提供了现成的属性变化委托。  
所以属性值变化时，可以及时触发更新。

\subsubsection{为什么 Tag / GE / Ability 用轮询刷新}
因为当前项目阶段里，还没有为这些内容单独封装完整的事件驱动同步机制，比如：

\begin{itemize}
	\item Tag 增删事件聚合
	\item GE 增删变化通知
	\item Ability 授予或移除时的 UI 同步
\end{itemize}

所以当前实现采用了一个很实用的工程折中：

\begin{itemize}
	\item 属性：用 Delegate
	\item 其他信息：面板可见时按固定间隔轮询刷新
\end{itemize}

\subsection{这套代码里最值得学习的 4 个设计点}

\subsubsection{1. HUD 持有调试 UI，而不是 Character 持有}
这是一个比较清晰的职责分层：

\begin{itemize}
	\item \texttt{Character} 管角色能力系统
	\item \texttt{HUD} 管显示
	\item \texttt{PlayerController} 管输入
\end{itemize}

这样结构更稳，也更方便以后维护。

\subsubsection{2. Widget 不自己乱找对象，而是显式传入依赖}
\texttt{InitializeWithAbilitySystem(...)} 很关键。  
它不是在 Widget 里到处 \texttt{GetPlayerCharacter()}、\texttt{GetPlayerState()}，而是由外部把：

\begin{itemize}
	\item \texttt{ASC}
	\item \texttt{AttributeSet}
\end{itemize}

明确传进来。

这是一种显式依赖注入的思路，通常比在内部随便查对象更清晰。

\subsubsection{3. 用“缓存 + 统一更新显示”，而不是边算边改 UI}
代码不是在每个地方都直接 \texttt{SetText(...)}，而是先把数据整理成缓存，比如：

\begin{itemize}
	\item \texttt{CachedHealth}
	\item \texttt{CachedOwnedTagsText}
	\item \texttt{CachedActiveEffectsText}
	\item \texttt{CachedAbilitiesText}
\end{itemize}

最后再统一调用 \texttt{UpdateDisplayedValues()}。

这种模式很适合做：

\begin{itemize}
	\item 调试面板
	\item 状态面板
	\item 角色信息面板
\end{itemize}

\subsubsection{4. 只有可见时才轮询}
在 \texttt{NativeTick(...)} 里，不是永远高频刷新，而是：

\begin{itemize}
	\item 面板可见时才刷新
	\item 并且有 \texttt{RefreshIntervalSeconds} 控制刷新间隔
\end{itemize}

这是一个很实用的 UI 性能优化习惯。

\subsection{建议你边读边回答的 8 个问题}
如果你想真正吃透这套代码，建议边读边回答下面这些问题：

\begin{enumerate}
	\item \texttt{F12} 按下后，完整调用链是怎样的？
	\item 调试面板是谁创建的？谁持有它？
	\item \texttt{ASC} 和 \texttt{AttributeSet} 是从哪里传进 Widget 的？
	\item 为什么属性用 Delegate，而其他信息用定时刷新？
	\item \texttt{RefreshActiveGameplayEffects()} 里的每一行信息分别来自哪个 GAS API？
	\item \texttt{GetActivatableAbilities()} 返回的是什么？
	\item \texttt{BuildWidgetTree()} 里 Widget 的层级结构是什么？
	\item 如果要加一个 \texttt{Armor} 文本，需要改哪几处？
\end{enumerate}

如果这 8 个问题都能回答出来，这套代码基本就算真正吃透了。

\subsection{建议你自己动手做的一个最小练习}
建议做一个小练习：加一行显示 \texttt{MaxHealth - Health}。

例如增加一个：

\begin{itemize}
	\item \texttt{MissingHealthText}
\end{itemize}

用于显示：

\begin{itemize}
	\item \texttt{Missing Health}
\end{itemize}

这个练习会让你完整走一遍扩展流程：

\begin{enumerate}
	\item 在 \texttt{.h} 里加一个 \texttt{UTextBlock*}
	\item 在 \texttt{BuildWidgetTree()} 里创建它
	\item 在 \texttt{UpdateDisplayedValues()} 里给它设置文本
\end{enumerate}

这是最适合入手的练习：改动小，但能把整体结构摸熟。
\chapter{Gameplay Effect}
\section{让返回值变成蓝的执行引脚}
\begin{code}
	UFUNCTION(BlueprintCallable, meta = (ExpandBoolAsExecs = "ReturnValue"))
	bool ApplyEffectToTarget(AActor* TargetActor, TSubclassOf<UGameplayEffect> InGameplayEffectClass);
\end{code}
如果是枚举类型，则可以：
\begin{code}
	meta = (ExpandEnumAsExecs = "SomeEnumParam")
\end{code}
\section{Effect Stacking}

定义\textbf{同一个} GameplayEffect 定义被\uwave{重复}应用时，\textbf{非瞬时} GE（non-instant GE）应当如何叠层。

也就是说，它解决的是这个问题：
当一个\emph{持续型 / 周期型 / 有时长}的 GameplayEffect，被再次应用时，是创建一个新实例，还是并到已有实例上形成 stack？

注意两个关键词：
\begin{enumerate}
	\item 同一个 GameplayEffect 定义
	\item 非 Instant GE
\end{enumerate}

所以它主要影响的是：
\begin{itemize}
	\item Duration
	\item Infinite
\end{itemize}
通常也包括带周期的持续效果。而 Instant GE 一般不会靠这个来“叠层存在”，因为它是瞬时执行完就没了

本节围绕 \texttt{Gameplay Effect} 的 \textbf{Stacking}（叠加）展开，重点说明了：
\begin{enumerate}
	\item \textbf{没有 Stacking} 并不等于“不能叠加”，而是“同一个效果被独立施加多次”；
	\item 真正的 \textbf{Stacking} 是把多次施加视为同一个可控的堆叠实例，并通过一组策略控制其
	\textbf{层数上限、持续时间、周期重置、过期行为}；
	\item \texttt{Aggregate by Source} 与 \texttt{Aggregate by Target} 的差异，本质上是
	``\textbf{层数限制按谁来统计}''。
\end{enumerate}

\subsection{没有 Stacking 时，发生了什么？}

当 \texttt{Stacking Type = None} 时，效果并不会进入“叠层管理”模式，而是每次拾取晶体都
\textbf{独立创建一个新的 Gameplay Effect 实例}。

例如：
\begin{itemize}
	\item 每个 Mana Crystal 持续 $1$ 秒；
	\item 周期（Period）为 $0.1$ 秒；
	\item 每个周期给角色 $+1$ Mana；
	\item 那么一个晶体理论上总共给出 $	10 \times 1 = 10 \text{ 点 Mana}	$
\end{itemize}

如果在很短时间内连续拾取 $3$ 个晶体，那么不是“一个效果有了三层”，而是：
\[
\text{独立效果}_1 + \text{独立效果}_2 + \text{独立效果}_3
\]
同时运行。于是 Mana 增长速度会明显变快，因为三个周期效果都在各自结算。

这正是视频里强调的关键点：“\textbf{看起来像 stacking，但本质上不是 stacking；只是同一个效果被施加了三次。}”

\subsection{真正的 Stacking 在解决什么问题？}

真正的 Stacking 解决的是：\textbf{当一个持续型效果尚未结束时，又来了同类效果，该如何处理？}

常见设计目标包括：
\begin{itemize}
	\item 不希望数值增长速度无限叠高；
	\item 但希望再次拾取时可以 \textbf{刷新持续时间}；
	\item 或者希望只增加层数，不增加频率；
	\item 或者希望超出上限时直接拒绝施加；
	\item 或者希望过期时一次性清空所有层；
	\item 或者希望过期时只减一层，并把持续时间重新开始计时。
\end{itemize}

也就是说，Stacking 的重点不是“能不能叠”，而是：	\textbf{叠了之后，系统应该按什么规则继续运行。}

\subsection{三种 Stacking Type 的意义}

\begin{longtable}{p{0.2\linewidth}p{0.25\linewidth}p{0.55\linewidth}}
	\toprule
	\textbf{类型} & \textbf{含义} & \textbf{效果}\\
	\midrule
	\endfirsthead
	\toprule
	\textbf{类型} & \textbf{含义} & \textbf{效果}\\
	\midrule
	\endhead
	\texttt{None} &
	不做堆叠聚合 &
	每次施加都是独立实例；持续时间、周期、结算彼此分开。\\
	\texttt{Aggregate by Source} &
	按施加者（Source）统计堆叠 &
	每个 Source 对同一 Target 都有自己独立的堆叠计数。\\
	\texttt{Aggregate by Target} &
	按受体（Target）统计堆叠 &
	无论来自哪个 Source，最终都并入 Target 的同一个堆叠计数。\\
	\bottomrule
\end{longtable}

\subsection{Aggregate by Source 与 Aggregate by Target 的核心区别}

假设 \texttt{Stack Limit Count = 2}。

\subsubsection{Aggregate by Source}

它表示：\textbf{每个 Source 最多只能给目标叠到 2 层。}

若 Source A 连续施加三次：
\[
A \to T: 1, 2, \cancel{3}
\]
第三次失败，因为 A 自己的层数已经达到上限。

但如果随后 Source B 再来一次：
\[
B \to T: 1
\]
这是允许的，因为 B 的“来源计数”还是第一次。

于是 Target 身上可能出现这样的状态：
\[
\text{A 提供 2 层} + \text{B 提供 1 层}
\]
虽然总层数看起来像 $3$，但它并没有违反规则，因为规则是“\textbf{每个 Source 最多 2 层}”。

\subsubsection{Aggregate by Target}

它表示：\textbf{Target 身上的总层数最多只能到 2 层，不管是谁施加的。}

如果 A 已经给 T 叠了两层：
\[
A \to T: 1, 2
\]
此时 B 再施加一次：
\[
B \to T: \cancel{3}
\]
依然失败，因为 Target 的总堆叠已经满了。

\subsubsection{一句话区分}

\begin{itemize}
	\item \texttt{Aggregate by Source}：限制是“\textbf{每个来源各算各的}”；
	\item \texttt{Aggregate by Target}：限制是“\textbf{目标总盘子只有这么大}”。
\end{itemize}

\subsection{字幕中的实验案例：为什么 Source 和 Target 在这里看起来差不多？}

视频里提到，在当前实验场景下，Mana Crystal 的效果是由场景中的 Actor 触发，但还没有明显体现出多个不同 Ability System Component 作为 Source 的复杂情况。因此在该演示里：

\begin{itemize}
	\item 选 \texttt{Aggregate by Source}；
	\item 或选 \texttt{Aggregate by Target}；
\end{itemize}

最终都表现为“只能叠到限定层数”。

这是因为当前案例并没有构造出多个不同 Source 同时给同一 Target 施加同一效果的复杂环境，所以两者的区别尚未完全显现。

\subsection{叠层相关的四个关键策略}

除了 \texttt{Stacking Type}，本节更重要的是四类控制参数。

\begin{longtable}{p{0.32\linewidth}|p{0.64\linewidth}}
	\toprule
	\textbf{参数} & \textbf{作用}\\
	\midrule
	\endfirsthead
	\toprule
	\textbf{参数} & \textbf{作用}\\
	\midrule
	\endhead
	\texttt{Stack Limit Count} &
	允许的最大层数。超过后，新效果不能再成功施加。\\
	\texttt{Stack Duration Refresh Policy} &
	当成功新增一层时，是否刷新整个效果的持续时间。\\
	\texttt{Stack Period Reset Policy} &
	当成功新增一层时，是否重置周期计时器（Period）。\\
	\texttt{Stack Expiration Policy} &
	当持续时间结束时，整个堆叠如何处理。\\
	\bottomrule
\end{longtable}

\subsection{Duration Refresh Policy：叠层时要不要刷新持续时间？}

字幕里重点对比了两种策略：
\begin{enumerate}
	\item \textbf{Never Refresh}
	
	持续时间一旦开始计时，就按原定时间结束，不因后续叠层而延长。例如：
	\begin{itemize}
		\item 第一个晶体在 $t=0$ 拾取，持续到 $t=1$；
		\item 第二个晶体在 $t=0.5$ 拾取；
		\item 如果策略是 \texttt{Never Refresh}，那么整个效果仍可能在原本截止点附近结束；
		\item 所以玩家无法拿到完整的“两秒收益”。
	\end{itemize}
	
	这正是视频中“没能涨到预期值”的原因：不是第二次拾取没生效，而是\textbf{总持续时间没有被刷新}。
	
	\item \textbf{Refresh on Successful Application}
	
	每当成功新增一层，就把持续时间重新开始计算。
	
	如果基础持续时间是 $1$ 秒，那么在效果快结束前又获得一层，就相当于又续了 $1$ 秒。  
	这更适合“不断拾取就不断续杯”的设计思路。
\end{enumerate}
\subsection{Period Reset Policy：叠层时要不要重置周期？}

假设周期为：
$
\text{Period} = 0.1\text{ 秒}
$

它控制的不是总持续时间，而是“\textbf{下一次周期触发什么时候到来}”。
\begin{enumerate}
	\item \textbf{Reset on Successful Application}
	
	新增一层后，周期重新开始计时。  
	也就是如果你原本只差一点就要触发下一次加 Mana，叠层之后会重新等一个完整周期。
	
	\item \textbf{Never Reset}
	
	新增一层不会打断当前周期节奏。  
	原本什么时候该结算，仍然什么时候结算。
\end{enumerate}
\subsubsection{直观理解}

\begin{itemize}
	\item \textbf{Duration} 关心的是：这个效果还剩多久结束；
	\item \textbf{Period} 关心的是：下一跳数值变化何时发生。
\end{itemize}

两者看起来接近，但控制的是不同维度。

\subsection{Expiration Policy：效果过期后如何处理？}

这是本节最容易混淆、也最有设计空间的部分。

\subsubsection{Clear Entire Stack}

一旦持续时间结束，整个堆叠直接清空。

若当前有 $3$ 层，那么到点后结果是：
\[
3 \to 0
\]

特点：
\begin{itemize}
	\item 行为简单；
	\item 容易理解；
	\item 适合“过期就整个结束”的增益或减益。
\end{itemize}

\subsubsection{Remove Single Stack and Refresh Duration}

每次过期时，只移除一层，并把持续时间重新开始。

若当前有 $3$ 层，那么可能表现为：
\[
3 \to 2 \to 1 \to 0
\]
而且每次层数减少后，都会重新开始一轮持续时间。

这相当于让堆叠逐层衰减，而不是瞬间消失。  
适合设计“层数会慢慢掉”的 Buff 或 Debuff。

\subsubsection{Refresh Duration}

字幕特别强调：这个选项可能导致一种“近似无限持续”的效果。

因为它的语义是：\textbf{当栈过期时，刷新持续时间}。  
如果连最后一层也这样处理，那么：
\[
1 \xrightarrow{\text{过期}} 1 \xrightarrow{\text{过期}} 1 \xrightarrow{\text{过期}} \cdots
\]

结果就是效果不断自我续期，看起来永远不会结束。  
视频里明确指出，这可用于通过回调（如 \texttt{onStackCountChange}）手动管理层数递减。

\subsection{用一个公式理解视频中的 Mana 例子}

设：
\[
d = 1\text{ 秒}, \qquad p = 0.1\text{ 秒}, \qquad \Delta m = 1
\]

则单层效果理论收益为：
\[
\frac{d}{p} \cdot \Delta m = \frac{1}{0.1}\cdot 1 = 10
\]

如果 \texttt{Stacking Type = None}，且短时间拾取 $n$ 个晶体，则近似同时存在 $n$ 个独立效果：
\[
\text{收益速率} \approx n \times \frac{\Delta m}{p}
\]
因此数值上涨会非常快。

如果使用真正的 stacking，且限制层数为 $1$，同时“成功施加时刷新持续时间”，那么再拾取晶体只是在做：
\[
\text{延长持续时间}
\]
而不是提升增长速率。  
所以你看到的是“持续更久地每 $0.1$ 秒加 $1$”，而不是“每 $0.1$ 秒加很多”。

\subsection{一个时间线示意图}

下面用一个简单的 TikZ 时间线表示“叠层但不提高速率，只刷新时长”的情形。

\begin{center}
	\begin{tikzpicture}[x=1cm,y=1cm,>=stealth]
		\draw[->] (0,0) -- (11,0) node[right] {$t$};
		\foreach \x in {0,1,...,10} {
			\draw (\x,0.1) -- (\x,-0.1) node[below=4pt] {\small \x};
		}
		
		\draw[thick,blue] (0,0.6) -- (4,0.6);
		\node[blue] at (2,0.95) {\footnotesize 第一次拾取后的持续时间};
		
		\draw[thick,blue] (3,1.3) -- (7,1.3);
		\node[blue] at (5,1.65) {\footnotesize 在 $t=3$ 再拾取，刷新持续时间};
		
		\draw[thick,blue] (6,2.0) -- (10,2.0);
		\node[blue] at (8,2.35) {\footnotesize 在 $t=6$ 再拾取，再次刷新};
		
		\filldraw[red] (0,0.6) circle (2pt);
		\filldraw[red] (3,1.3) circle (2pt);
		\filldraw[red] (6,2.0) circle (2pt);
		
		\node[red] at (0,0.25) {\footnotesize 拾取};
		\node[red] at (3,1.0) {\footnotesize 拾取};
		\node[red] at (6,1.7) {\footnotesize 拾取};
	\end{tikzpicture}
\end{center}

这个图表达的是：
\begin{itemize}
	\item 玩家多次拾取并没有让“每次加 Mana 的速度”越来越快；
	\item 而是让同一个效果的有效持续区间不断向后延长。
\end{itemize}

\subsection{本节字幕最终推荐的配置思路}

视频后半段实际上给出了一个比较实用的设计方向，尤其适合拾取类资源晶体：

\begin{itemize}
	\item \texttt{Stacking Type = Aggregate by Target}
	\item \texttt{Stack Limit Count = 1} 或较小值
	\item \texttt{Stack Duration Refresh Policy = Refresh on Successful Application}
	\item \texttt{Stack Period Reset Policy = Reset on Successful Application}
	\item \texttt{Stack Expiration Policy = Clear Entire Stack} 或 \texttt{Remove Single Stack and Refresh Duration}
\end{itemize}

这样做的设计收益是：
\begin{itemize}
	\item 连续拾取时不会导致属性暴涨失控；
	\item 但玩家会感觉“再次拾取是有价值的”，因为效果可以续时；
	\item 系统行为稳定，可预测，便于平衡；
	\item 能对“瞬间吃到太多晶体”的情况做一定节流（throttling）。
\end{itemize}

\subsection{Crystal Heal 的设置意图}

字幕最后把 \texttt{Crystal Heal} 也配置成类似逻辑，其核心目的不是“让治疗更强”，而是：

\begin{itemize}
	\item 限制同类拾取物在极短时间内造成过强的瞬时收益；
	\item 给拾取效果一个可控、可平衡的叠层规则；
	\item 避免玩家一次性卷入多个晶体时造成难以调参的爆发恢复。
\end{itemize}

\subsection{本节真正应该记住的结论}

\begin{enumerate}
	\item \textbf{None} 不代表没有叠加现象，而是“多次独立施加”。
	\item 真正的 \textbf{Stacking} 是把多次施加纳入统一的堆叠管理体系。
	\item \texttt{Aggregate by Source} 与 \texttt{Aggregate by Target} 的区别，在于层数上限究竟按 Source 统计还是按 Target 统计。
	\item \texttt{Stack Limit Count} 决定最多允许几层。
	\item \texttt{Duration Refresh Policy} 决定新层加入时是否续时。
	\item \texttt{Period Reset Policy} 决定新层加入时是否重置下一次周期结算。
	\item \texttt{Expiration Policy} 决定过期时是整栈清空、逐层衰减，还是被刷新成近似无限效果。
	\item 对拾取类持续效果来说，Stacking 的核心价值往往不是“让数值更猛”，而是“让行为更可控、更容易平衡”。
\end{enumerate}

\section{Copilot：对EGameplayEffectStackingType 解释}

我先直接解释这个枚举的语义，再结合 GAS 的实际行为帮你建立直觉。

\begin{itemize}
	\item 解释 \texttt{EGameplayEffectStackingType} 是干什么的
	\item 逐个说明 3 个枚举值的含义
	\item 说明“谁和谁会叠层”
	\item 给出具体例子
	\item 说明它和“普通重复应用但不叠层”的区别
\end{itemize}

\subsection{它是什么}

\texttt{EGameplayEffectStackingType} 定义在：

\begin{itemize}
	\item \texttt{GameplayAbilities} 插件的 \texttt{GameplayEffectTypes.h}
\end{itemize}

它的注释原文意思是：

\begin{quote}
	定义同一个 GameplayEffect 定义被重复应用时，非瞬时 GE（non-instant GE）应当如何叠层。
\end{quote}

也就是说，它解决的是这个问题：

\begin{quote}
	当一个\dotuline{持续型 / 周期型 / 有时长}的 \texttt{GameplayEffect} 被再次应用时，  
	是创建一个新实例，还是并到已有实例上形成 stack？
\end{quote}

注意两个关键词：

\begin{enumerate}
	\item 同一个 GameplayEffect 定义
	\item 非 Instant GE
\end{enumerate}

所以它主要影响的是：

\begin{itemize}
	\item \texttt{Duration}
	\item \texttt{Infinite}
	\item 通常也包括带周期的持续效果
\end{itemize}

而 \texttt{Instant GE} 一般不会靠这个来“叠层存在”，因为它是瞬时执行完就没了。

\subsection{枚举定义}


\begin{code}
	/** Defines how multiple instances of the same GameplayEffect definition stack when applied (non-instant GE's). */
	UENUM()
	enum class EGameplayEffectStackingType : uint8
	{
		/** No stacking. Each application is a separate instance. */
		None UMETA(DisplayName = "No Stacking"),
		
		/** Stacks only when the same source reapplies the effect. */
		AggregateBySource UMETA(DisplayName = "Stack Per Source"),
		
		/** Stacks only when the effect is reapplied to the same target. */
		AggregateByTarget UMETA(DisplayName = "Stack Per Target"),
	};
\end{code}

\subsection{\texttt{None}}

\subsubsection*{含义}
不叠层。每次应用都是一个独立实例。

即使是同一个 GE，被重复施加，也不会合并成一个 stack，而是会出现多个独立 Active Gameplay Effect。

\subsubsection*{直觉理解}
你可以把它理解成：

\begin{itemize}
	\item 第一次上 Buff：生成一条效果
	\item 第二次再上同一个 Buff：再生成一条效果
	\item 两条效果各自有自己的：
	\begin{itemize}
		\item 持续时间
		\item 周期 Tick
		\item 移除时机
	\end{itemize}
\end{itemize}

\subsubsection*{举例}
比如有个持续 5 秒的中毒效果 \texttt{GE\_Poison}：

\begin{itemize}
	\item 第 0 秒上一次
	\item 第 2 秒再上一遍
\end{itemize}

如果 stacking type 是 \texttt{None}，那通常会变成：

\begin{itemize}
	\item 第一条毒：0\textasciitilde5 秒
	\item 第二条毒：2\textasciitilde7 秒
\end{itemize}

它们是两条独立效果，不是“2 层毒”。

\subsubsection*{适用场景}
适合这些情况：

\begin{itemize}
	\item 你希望每次施加都是独立实例
	\item 每个来源的效果都要完整保留
	\item 例如多个投射物命中后，各自产生独立 DOT
\end{itemize}
\textbf{副作用：}由于重叠，毒效果在$2\to 5$时会加速。
\subsection{\texttt{AggregateBySource}}

\subsubsection*{含义}
按 Source 聚合叠层。只有“同一个来源”再次施加时，才会叠到一起。

这里的 source 可以直观理解为：

\begin{itemize}
	\item 同一个施法者
	\item 同一个来源 \textbf{ASC}
	\item 同一个效果施加来源
\end{itemize}

\subsubsection*{直觉理解}
规则可以简化成：

\begin{itemize}
	\item 同一个敌人给你重复上同一个 GE：会叠层
	\item 不同敌人给你上同一个 GE：不会叠在同一堆里，而是分开算
\end{itemize}

\subsubsection*{举例}
有两个敌人：

\begin{itemize}
	\item 敌人 A
	\item 敌人 B
\end{itemize}

都能给玩家施加 \texttt{GE\_Burn}

如果 stacking type 是 \texttt{AggregateBySource}：

\begin{itemize}
	\item A 连续打你 3 次 → 形成 A 的 3 层 Burn
	\item B 连续打你 2 次 → 形成 B 的 2 层 Burn
\end{itemize}

结果通常不是“总共 5 层一个 Burn”，而是更像：

\begin{itemize}
	\item 一组来自 A 的 stack
	\item 一组来自 B 的 stack
\end{itemize}

\subsubsection*{适用场景}
适合这些设计：

\begin{itemize}
	\item 你想区分不同施法者的贡献
	\item 每个来源各自维护自己的层数/持续时间
	\item MMO、ARPG 里常见的“每个施法者单独挂 DOT”
\end{itemize}

比如：

\begin{itemize}
	\item 玩家 A 的流血
	\item 玩家 B 的流血
\end{itemize}

虽然是同一个 GE 定义，但你不希望他们互相挤掉或合并成一组。

\subsection{\texttt{AggregateByTarget}}

\subsubsection*{含义}
按 Target 聚合叠层。只要是对同一个目标重复施加这个 GE，就会并到同一组 stack。

不太关心是谁施加的，重点是：

\begin{itemize}
	\item 目标相同
	\item GE 定义相同
\end{itemize}

那就聚合。

\subsubsection*{直觉理解}
规则可以近似理解成：

\begin{itemize}
	\item 谁来施加不重要
	\item 只要大家都往同一个目标身上挂这个 GE，就叠同一组
\end{itemize}

\subsubsection*{举例}
还是两个敌人：

\begin{itemize}
	\item 敌人 A
	\item 敌人 B
\end{itemize}

都给玩家施加 \texttt{GE\_Burn}

如果 stacking type 是 \texttt{AggregateByTarget}：

\begin{itemize}
	\item A 给玩家 2 次
	\item B 再给玩家 3 次
\end{itemize}

最后在玩家身上，通常表现成：

\begin{itemize}
	\item 一组总共 5 层的 Burn
\end{itemize}

而不是分 A 一组、B 一组。

\subsubsection*{适用场景}
适合这些设计：

\begin{itemize}
	\item 你不关心是谁叠的
	\item 你只关心目标身上当前总共有几层
	\item 比如“易伤最多 5 层”、“灼烧最多 10 层”
\end{itemize}

这种玩法设计更偏“目标状态统一管理”。

\subsection{一张表看懂}

\begin{longtable}{|p{0.22\textwidth}|p{0.28\textwidth}|p{0.38\textwidth}|}
	\hline
	\textbf{枚举值} & \textbf{含义} & \textbf{谁会叠在一起} \\
	\hline
	\endfirsthead
	
	\hline
	\textbf{枚举值} & \textbf{含义} & \textbf{谁会叠在一起} \\
	\hline
	\endhead
	
	\hline
	\endfoot
	
	\hline
	\endlastfoot
	
	\texttt{None} & 不叠层 & 谁都不叠，每次都是独立实例 \\
	\hline
	\texttt{AggregateBySource} & 按来源叠层 & 同一个来源对同一个目标重复施加时叠层 \\
	\hline
	\texttt{AggregateByTarget} & 按目标叠层 & 只要对同一个目标施加同一个 GE，就叠到一起 \\
	\hline
\end{longtable}

\subsection{它到底在决定什么}

它决定的不是“数值怎么算”本身，而是更前一步：

\begin{quote}
	多个应用记录是否合并为同一个 Active Gameplay Effect Stack
\end{quote}

一旦合并成 stack，后续很多行为都会受影响，比如：
\begin{enumerate}
	\item \texttt{StackCount}
	\item 持续时间是否刷新
	\item 周期是否如何处理
	\item 到上限后怎么办
	\item 移除一层还是整个移除
\end{enumerate}

这些又会受其他 stacking 相关配置影响，比如：
\begin{enumerate}
	\item Stack Limit Count
	\item Duration Refresh Policy
	\item Period Reset Policy
	\item Expiration Policy
\end{enumerate}

所以你可以把 \texttt{EGameplayEffectStackingType} 看成：	\textbf{先决定按什么维度分组}，而不是：最终叠层时长和数值全部规则

\subsection{它和“重复应用但不叠层”的区别}

很多人第一次学 GAS 容易把这两个混掉。

\subsubsection*{情况 A：\texttt{None}}
表示：

\begin{itemize}
	\item 不做 stack 聚合
	\item 每次是独立 Active GE
\end{itemize}

\subsubsection*{情况 B：有 stacking type，但别的规则限制住了}
比如：

\begin{itemize}
	\item 虽然允许叠层
	\item 但 stack limit = 1
	\item 或应用时选择刷新/覆盖策略
\end{itemize}

这时表现可能看起来像“没叠起来”，但底层逻辑不一样。

所以要区分：

\begin{itemize}
	\item \texttt{None}：根本不进入“聚合叠层”路线
	\item \texttt{Aggregate...}：进入聚合路线，只是结果可能受 stack 上限和刷新策略约束
\end{itemize}

\subsection{一个非常实用的记忆法}

你可以这样记：

\begin{itemize}
	\item \texttt{None}：一发一个实例
	\item \texttt{AggregateBySource}：每个施法者各算各的层数
	\item \texttt{AggregateByTarget}：所有人往同一目标上叠同一层数池
\end{itemize}

\subsection{结合 ARPG / GAS 时的典型例子}

\subsubsection*{例子 1：中毒 DOT}
如果你想要：

\begin{itemize}
	\item 每个怪自己上的毒单独存在
\end{itemize}

那通常选：

\begin{itemize}
	\item \texttt{AggregateBySource}
\end{itemize}

因为 A 怪和 B 怪的毒不该混成一堆。

\subsubsection*{例子 2：易伤 Debuff}
如果你想要：

\begin{itemize}
	\item 目标身上的“易伤”就是统一最多 5 层
	\item 不管是谁打出来的
\end{itemize}

那通常选：

\begin{itemize}
	\item \texttt{AggregateByTarget}
\end{itemize}

\subsubsection*{例子 3：独立护盾效果}
如果你想要：

\begin{itemize}
	\item 每次加护盾都是一个独立计时器
	\item 互不合并
\end{itemize}

那通常选：

\begin{itemize}
	\item \texttt{None}
\end{itemize}

\subsection{再补一个边界点：它只针对“同一个 GE 定义”}

这是很重要的。

比如你有两个效果：

\begin{itemize}
	\item \texttt{GE\_Burn\_Small}
	\item \texttt{GE\_Burn\_Big}
\end{itemize}

即使它们看起来都叫“燃烧”，\emph{只要它们是两个不同的 GameplayEffect 类/定义，通常也不会因为这个枚举直接彼此叠在一起。}

所以这个 stacking type 讨论的是：	\textbf{同一个 GameplayEffect Definition 被重复应用时怎么办？} 而不是“不同效果之间怎么合并”。

\subsection{如果你只想记一句话}

\texttt{EGameplayEffectStackingType} 用来决定：同一个非瞬时 GameplayEffect 被重复施加时，是不叠层，还是按施法者分组叠层，还是按目标统一叠层。


% 需要的宏包：
% \usepackage{longtable,booktabs,array}
% \usepackage{tikz}
% \usetikzlibrary{arrows.meta,positioning}

\subsection{\texttt{EGameplayEffectStackingDurationPolicy}}

\textbf{作用：}
用于规定\emph{同一个 Gameplay Effect 发生堆叠时}，已有效果的\textbf{持续时间（Duration）}应如何处理。  
它关注的是：\textbf{新的一层叠加进来后，剩余时间怎么变。}

\vspace{0.5em}
\textbf{核心问题：}
\[
\text{新 Stack 成功应用} \;\Longrightarrow\; \text{Remaining Time 如何变化？}
\]

\begin{longtable}{@{}p{0.28\textwidth}p{0.32\textwidth}p{0.32\textwidth}@{}}
	\toprule
	\textbf{枚举值} & \textbf{含义} & \textbf{可直接记忆为} \\
	\midrule
	\endhead
	
	\texttt{RefreshOnSuccessfulApplication}
	&
	每次\textbf{成功叠层}时，都把持续时间\textbf{刷新}一次。
	&
	\textbf{重新计时} \\
	
	\texttt{NeverRefresh}
	&
	叠层只增加层数，\textbf{不刷新}当前剩余时间。
	&
	\textbf{层数变，时间不变} \\
	
	\texttt{ExtendDuration}
	&
	新叠上的这一层，会把它自己的持续时间\textbf{累加}到当前剩余时间上。
	&
	\textbf{剩余时间往后加} \\
	
	\bottomrule
\end{longtable}

\paragraph{直观例子}
假设某个可堆叠 GE 的基础持续时间是 $5\text{s}$。  
当前目标身上还剩 $2\text{s}$ 时，又成功叠加了一层：

\begin{itemize}
	\item \texttt{RefreshOnSuccessfulApplication}：
	$
	\text{Remaining Time} \gets 5\text{s}
	$\\
	即：\textbf{回满到新的完整时长}。
	
	\item \texttt{NeverRefresh}：
	$
	\text{Remaining Time} \gets 2\text{s}
	$\\
	即：\textbf{时间不动，只是层数增加}。
	
	\item \texttt{ExtendDuration}：
	$
	\text{Remaining Time} \gets 2\text{s} + 5\text{s} = 7\text{s}
	$\\
	即：\textbf{把新层的时长接到后面}。
\end{itemize}

\begin{figure}[H]
	\centering
	\begin{tikzpicture}[
		node distance=8mm and 10mm,
		box/.style={draw, rounded corners, align=center, minimum width=34mm, minimum height=8mm},
		arrow/.style={-Latex, thick}
		]
		\node[box] (start) {成功应用一个新 Stack};
		\node[box, below left=of start] (refresh) {\texttt{RefreshOnSuccessfulApplication}\\剩余时间重置为新时长};
		\node[box, below=of start] (never) {\texttt{NeverRefresh}\\剩余时间不变};
		\node[box, below right=of start] (extend) {\texttt{ExtendDuration}\\剩余时间 += 新时长};
		
		\draw[arrow] (start) -- (refresh);
		\draw[arrow] (start) -- (never);
		\draw[arrow] (start) -- (extend);
	\end{tikzpicture}
	\caption{\footnotesize\textbf{一图记忆}:\texttt{RefreshOnSuccessfulApplication}：\textbf{重置倒计时}； 		\texttt{NeverRefresh}：\textbf{只叠层，不续杯}；		\texttt{ExtendDuration}：\textbf{时间往后延长}。}
	\label{figRefresh}
\end{figure}

\paragraph{一句话理解}
\[
\boxed{
	\texttt{Refresh} = 重新开始
	\quad
	\texttt{NeverRefresh} = 不改时间
	\quad
	\texttt{Extend} = 累加时间
}
\]

% =========================
% GAS Stacking Notes
% =========================
% 建议宏包：
% \usepackage{longtable,booktabs,array}
% \usepackage{tikz}
% \usetikzlibrary{arrows.meta,positioning}
% \usepackage{amsmath}
% \usepackage{xcolor}